\subsection{Load balancing}

\begin{paragrafo}{Problema:}
    \begin{paragrafo}{Input:}
        Ho in input un numero $n \in \mathbb{N}^+$ di task e un numero $m \in \mathbb{N}^+$ di macchinari.

        Ognuna di queste task ha un tempo di esecuzione $t_0,t_1,\dots t_{n-1} \in \mathbb{N}^+$

        Indichiamo con $L_i$ il carico associato a una macchina $i$.

        Indichiamo con $L$ il costo, definito come  $L = \max_i{L_i}$
    \end{paragrafo}

    \begin{paragrafo}{Soluzioni ammesse:}
        Le soluzioni ammesse per questo problema sono delle partizioni di $n$ per $m$:

        \begin{center}
            $A_0,A_1,\dots ,A_{m-1}$
        \end{center}

        Il problema è di ricerca del minimo.
    \end{paragrafo}

    \begin{paragrafo}{Esempio grafico:}
        Supponiamo di avere 8 task, e 3 macchinari su cui eseguirle.

        \textbf{Tasks:}
        \begin{center}
            \memoria{3,2,3,1,1,4,5,1}
        \end{center}
        \smallskip
        Possiamo trovare una possibile partizione delle task per i macchinari del tipo:

        \smallskip
        \begin{center}
            \begin{tabular}{r@{\qquad}l@{\qquad}r}
                $m$ & tasks & $L_i$ \\
                0 & \memoriaunita{3,1,1,2} & 7 \\
                1 & \memoriaunita{5,1,3}   & 8 \\
                2 & \memoriaunita{4,1}     & 5
            \end{tabular}
        \end{center}
        \smallskip
        Vediamo facilmente che il costo $L$ è il massimo tra i 3, ovvero $8$.
        Possiamo allo stesso tempo affermare che non è la soluzione ottima, infatti, è semplice evidenziarne una migliore.

        \smallskip
        \begin{center}
            \begin{tabular}{r@{\qquad}l@{\qquad}r}
                $m$ & tasks & $L_i$ \\
                0 & \memoriaunita{5,1,1} & 7 \\
                1 & \memoriaunita{4,3}   & 7\\
                2 & \memoriaunita{3,2,1} & 6
            \end{tabular}
        \end{center}
        \smallskip
    \end{paragrafo}

    \newpage

    \begin{teorema}
        Il problema Load Balancing è NPO Completo

        \begin{rientro}
            Questo ci porta subito a studiare le soluzioni approssimate.
        \end{rientro}
    \end{teorema}

    Inserire algoritmo qui che non ho sbatti, e anche i vari tempi e cazzate varie.

    \begin{teorema}
        Greedy Load Balancing è 2 approssimato per Load Balancing.
    \end{teorema}

    \begin{dimostrazione}
        \osservazione{1}{$L^* \ge \frac{1}{m}\sum_jt_j$}

        \begin{giustificazione}{Dim 1.}{$\sum_{i} L_i^* = \sum_{j} t_j$}
            Per pidgeonholing, non è possibile che
            $\forall i L^*_i < \frac 1 m \sum_{j}t_j$

            quindi

            $\exists i L^*_i \ge \frac 1 m \sum_j t_j$
        \end{giustificazione}

        \osservazione{2}{$L^* \ge \max_j t_j$}

        \separapassaggi

        Sia $\hat i$ la macchina più carica alla fine dell'esecuzione e $\hat j$ l'ultima task assegnata a $\hat i$.

        Allora possiamo dire che:

        \[
            L_{\hat i} - t_{\hat j} \le L_i'\ \forall i \in m \le L_{i}
        \]

        A questo punto sommiamo le disequazioni:

        \[
            L_{\hat i} - t_{\hat j} \le L_i'
        \]
        \[
            m(L_{\hat i} - t_{\hat j}) \le \sum_i L_i
        \]
        \[
            L_{\hat i} - t_{\hat j} \le \frac 1 m \sum_i L_i
            = \frac 1 m \sum_j t_j\le L^*
        \]
        \[
            L = L_{\hat i} = (L_{\hat i} - t_{\hat j}) + t_{\hat j}
            \le L^* + t_{\hat j}\le L^*+L^*
        \]
        \[
            L\le2L^*
        \]

        \corollario{Load Balancing $\in$ 2-APX}
    \end{dimostrazione}

    \begin{teorema}
        $\forall \epsilon > 0$ esiste un input in cui Greedy Load Balancing produce una soluzione di costo:
        \[
            2-\epsilon \le \frac{L}{L^*} \le 2
        \]
    \end{teorema}

    \begin{dimostrazione}
        Prendiamo $m > \frac 1 \epsilon, \text{e } n = (m-1)+1$, con i primi $m-1$ input pari a $1$ el'ultimo pari a $m$.

        Applicando l'algoritmo su questi input otteniamo:
        \begin{center}
            \begin{tabular}{r@{\qquad}l@{\qquad}r}
                $m$ & tasks & $L_i$ \\
                0 & \memoriaunita{1,1,1} & 13 \\
                1 & \memoriaunita{1,1,3} & 3 \\
                2 & \memoriaunita{1,1}   & 3
            \end{tabular}
        \end{center}
        \[
            L = 2m-1
        \]
        Notiamo invece il caso ottimale su questo input:
        \begin{center}
            \begin{tabular}{r@{\qquad}l@{\qquad}r}
                $m$ & tasks & $L_i$ \\
                0 & \memoriaunita{1,1,1} & 3 \\
                1 & \memoriaunita{1,1,1} & 3 \\
                2 & \memoriaunita{1,3}   & 4
            \end{tabular}
        \end{center}
        \[
            L = 4, \text{ ottimo}
        \]
        \begin{digressione}
            L'input definito è inoltre uno dei pochi casi in cui la media equivale al caso ottimo.
            \[
                \frac{m(m-1)+m}{m}= m-1+1 = \lceil m \rceil
            \]
        \end{digressione}

        \[
            \frac L {L^*} = \frac {2m-1} {m}
            = 2 - \frac 1 m \ge 2-\epsilon
        \]
    \end{dimostrazione}

    Se ordino Greedy Load Balancing, perdo la sua proprietà di essere ONLINE, ma ne guadagno in approssimazione, infatti:

    \begin{lstlisting}[style=mystyle,mathescape=true,caption=Offline Greedy Load Balancing,label={alg:offline-greedy}]
    Sorto i task in ordine non crescente di durata
    Applico Greedy Load Balancing
\end{lstlisting}

    \begin{teorema}
        Offline Greedy Load Balancing è $\frac{3}{2}$ approssimato per Load Balancing.
    \end{teorema}

    \begin{dimostrazione}
        \osservazione{1}{Se $n \le m$, troviamo l'ottimo}

        \begin{rientro}
            Assumiamo quindi $n > m$
        \end{rientro}

        \osservazione{2}{$L^* \ge 2t_m$ dato che}
            \begin{giustificazione}{Dim 2.}{}
                $\underbrace{t_0 \ge t_1 \ge \dots \ge t_{m-1}}_{m\text{ task}}
                \ge t_m \dots \text{ sempre per pidgeonholing}$

                $\exists i\ L^* \ge L^*_i \ge t_{1} + t_{2} \ge t_m + t_m$

                quindi

                $L^* \ge 2t_m$
            \end{giustificazione}

            Sia $\hat i$ la macchina più carica alla fine dell'esecuzione.

            Possiamo assumere che ad essa siano assegnati almeno due task, altrimenti avremmo l'ottimo.

            Sia $\hat j$ l'ultimo task assegnato:

            \[
                \hat j \ge m
            \]

            quindi

            \[
                t_{\hat j} \le t_m \le \frac 1 2 L^*
            \]

            quindi

            \[
                L = L_{\hat i} = L_{\hat i} - t_{\hat j} + t_{\hat j}
                \le \frac 3 2 L^*
            \]

            \begin{nota}{1}
                L'analisi è migliorabile, dimostrando che l'algoritmo è
                $\frac 4 3 $ approssimato.
                \fonte{Graham 1969}
            \end{nota}

            \begin{nota}{2}
                Load Balancing $\notin $ FPTAS (assunto $P \not = NP$)
                \fonte{Hochbaum-Shmops 1988}
            \end{nota}
    \end{dimostrazione}
\end{paragrafo}
